\documentclass[a4paper]{article}

% Basic packages
% Español (México) con XeLaTeX: fontspec para fuentes Unicode, polyglossia
% para la división silábica y los nombres del documento en la variante mexicana.
\usepackage{fontspec}
\usepackage{polyglossia}
\setdefaultlanguage[variant=mexican]{spanish}
% Los nombres chinos van como «pinyin (original)»: xeCJK da a los caracteres
% Han su propia fuente; sin ella XeLaTeX los omite con un aviso.
% FandolSong no cubre algunos caracteres de nombres (U+73FA); WenQuanYi Zen Hei sí.
\usepackage[AutoFallBack=true]{xeCJK}
\setCJKmainfont{FandolSong-Regular.otf}
\setCJKfallbackfamilyfont{\CJKrmdefault}{WenQuanYi Zen Hei}
% polyglossia no traduce los nombres de listings.
\AtBeginDocument{\providecommand{\lstlistingname}{}\renewcommand{\lstlistingname}{Listado}%
  \providecommand{\lstlistlistingname}{}\renewcommand{\lstlistlistingname}{Índice de listados}}
\usepackage{amsmath, amssymb}
\usepackage{graphicx}
\usepackage[margin=2.5cm]{geometry}
\usepackage[most]{tcolorbox}
\usepackage{etoolbox}
\usepackage{listings}
\usepackage{hyperref}
\usepackage{booktabs}
\usepackage{subcaption}
\usepackage{float}
\usepackage{tikz}
\IfFileExists{pgfplots.sty}{
    \usepackage{pgfplots}
    \pgfplotsset{compat=1.18}
}{}

% Highlight boxes
\newtcolorbox{knowledgebox}[1]{
    enhanced, colback=blue!5!white, colframe=blue!75!black, colbacktitle=blue!75!black,
    coltitle=white, fonttitle=\bfseries, title={#1}, attach boxed title to top left={yshift=-2mm, xshift=2mm},
    boxrule=1pt, sharp corners
}

\newtcolorbox{importantbox}[1]{
    enhanced, colback=yellow!10!white, colframe=yellow!80!black, colbacktitle=yellow!80!black,
    coltitle=black, fonttitle=\bfseries, title={#1}, sharp corners
}

\newtcolorbox{warningbox}[1]{
    enhanced, colback=red!5!white, colframe=red!75!black, colbacktitle=red!75!black,
    coltitle=white, fonttitle=\bfseries, title={#1}, sharp corners
}

% Listings
\lstset{
    language=Python,
    basicstyle=\ttfamily\small,
    keywordstyle=\color{blue},
    stringstyle=\color{red!60!black},
    commentstyle=\color{green!60!black},
    breaklines=true,
    frame=single,
    numbers=left,
    numberstyle=\tiny\color{gray},
    captionpos=b,
    extendedchars=true
}

% Front-page metadata
\newcommand{\notetitle}{[LLM Agents SP25] Learning to Reason with LLMs — Jason Weston}
\newcommand{\noteauthors}{Elaboradas a partir de la clase de Jason Weston}
\newcommand{\notedate}{2025-02-10}
\newcommand{\videochannel}{Berkeley RDI}
\newcommand{\videopublishdate}{2025-02-10}
\newcommand{\videoduration}{1:16:47}
\newcommand{\videourl}{https://www.youtube.com/watch?v=_MNlLhU33H0}
\newcommand{\videocoverpath}{cover.jpg}
\newcommand{\repourl}{https://github.com/hqhq1025/ai-course-notes}

% Footer with repo info
\usepackage{fancyhdr}
\pagestyle{fancy}
\fancyhf{}
\fancyfoot[C]{\thepage}
\fancyfoot[R]{\footnotesize\color{gray}\href{\repourl}{AI Course Notes}}
\renewcommand{\headrulewidth}{0pt}
\renewcommand{\footrulewidth}{0pt}

\begin{document}

\begin{titlepage}
\centering
{\Large Notas de la entrevista\par}
\vspace{1.2cm}
{\huge\bfseries \notetitle\par}
\vspace{0.8cm}
{\large \noteauthors\par}
\vspace{0.3cm}
{\large \notedate\par}
\vspace{1.2cm}

\ifdefempty{\videocoverpath}{
{\small Al generar las notas, indique la ruta local de la portada del video.\par}
}{
\includegraphics[width=0.82\textwidth,height=0.45\textheight,keepaspectratio]{\videocoverpath}\par
}

\vfill
\begin{tcolorbox}[width=0.9\textwidth, colback=black!2!white, colframe=black!60, sharp corners]
\textbf{Autor o canal del video}: \videochannel\par
\textbf{Fecha de publicación}: \videopublishdate\par
\textbf{Duración del video}: \videoduration\par
\ifdefempty{\videourl}{}{
\textbf{Enlace del video}: \href{\videourl}{\nolinkurl{\videourl}}\par
}
\textbf{Página del proyecto}: \href{\repourl}{\nolinkurl{\repourl}}\par
\end{tcolorbox}
\end{titlepage}

\tableofcontents
\newpage

\section{Antecedentes y objetivos de la charla}
\subsection{El ponente y su lugar en el curso}
Esta charla pertenece al curso avanzado UC Berkeley CS294-280 Sp25 y la dirige Jason Weston, de Meta Research. Desde el inicio, la charla señala que la investigación actual en reasoning atraviesa saltos comparables a o1/R1 y DeepSeek-R1: los modelos avanzan con rapidez en los benchmark multimodales, pero cada salto también revela que los errores de razonamiento no se detectan a tiempo. Por eso, esta sección busca reconstruir el lugar de la charla: es a la vez un resumen de la historia reciente de Chain-of-Thought (que en realidad se popularizó hasta 2022) y un puente entre «dejar que el modelo defina por sí mismo la tarea de reasoning» y «dejar que él mismo juzgue el resultado».

\subsection{Propósito de la charla}
Weston resume el tema como self-improving language models: modelos que, durante el entrenamiento, aprenden poco a poco a formular sus propias preguntas, calificarse a sí mismos y optimizarse. Al inicio, el ponente menciona varias veces que «me parece que esta generación de investigadores de IA tiene mucha suerte de poder presenciar estos avances tan veloces» y subraya que «de forma análoga a un evento de benchmark como DeepSeek-R1, incluso en unos pocos meses se puede observar un cambio cualitativo en la reasoning accuracy». Por eso, a continuación las notas se desarrollan en torno a cuatro dimensiones: el objetivo del reasoning, el ciclo cerrado de entrenamiento, el sistema de evaluación y los retos del despliegue.

\begin{knowledgebox}{El contexto histórico de Chain-of-Thought}
El marco de entrenamiento de Chain-of-Thought se popularizó de manera explosiva a partir de 2022; Weston lo considera «un punto de partida rápido para el reasoning», pero también advierte que la verdadera capacidad de razonamiento exige que el modelo no solo produzca la cadena de razonamiento, sino que también sea capaz de evaluar cuál de las cadenas es más confiable.
\end{knowledgebox}

\subsection{Resumen de la sección}
Esta sección, a partir de hitos históricos (el auge de Chain-of-Thought, DeepSeek-R1, o1/R1) y de las reflexiones del ponente sobre el «self-improvement», establece el objetivo general y el criterio de evaluación de toda la charla. La pregunta central que hay que atender a continuación es: ¿cómo descomponer la «capacidad de introspección del modelo» en pasos concretos de entrenamiento?
\section{Rutas de razonamiento para la automejora}
\subsection{De Chain-of-Thought a la autorrecompensa}
Weston subraya que el crecimiento de la capacidad de razonamiento no puede depender solo de un prompt fuerte, sino que requiere que el propio modelo, durante la etapa de entrenamiento, «piense sobre su propio pensamiento». De ahí surgen los Self-Rewarding language models: durante el entrenamiento, el modelo desempeña a la vez el papel de generador de problemas, de quien los resuelve y de evaluador, y sustituye la calificación humana por una autoevaluación, de manera similar a como el mecanismo de ``verified response'' del Self-Feeding Chatbot se extiende a las tareas de razonamiento. Este proceso busca eliminar la confusión entre «memoria» y «razonamiento», por lo que los base data siguen incluyendo ejemplos que fuerzan el chain-of-thought, de modo que el reward model pueda identificar qué parte del reasoning explica la respuesta.

\subsection{El equilibrio entre el razonamiento y la autoevaluación}
Durante las iteraciones de self-improvement, el modelo cae fácilmente en un razonamiento encadenado excesivo: para obtener un reward más alto, genera chains cada vez más largas pero no necesariamente más correctas. Weston introduce el contraste ``self-rewarding vs self-feeding'' y señala que este último depende de la retroalimentación conversacional, mientras que el primero depende de la autoverificación (self-verification). Propone insertar ``cross checking'' en el entrenamiento self-rewarding: si el modelo no logra rechazar un error en el draft inicial, se le hace ``preguntarse a sí mismo'' una vez más.

\begin{warningbox}{Evitar el overthinking}
Por razones de seguridad, es necesario limitar la expansión del Chain-of-Thought, pues de lo contrario el inference compute se desperdicia en pensamientos repetidos. La introducción del Meta judge responde en parte a la necesidad de detectar a tiempo los casos en que ``incluso con una cadena de razonamiento larga, la respuesta sigue siendo incorrecta''.
\end{warningbox}

\subsection{Señales de entrenamiento y modelos de recompensa}
En la clase, Weston usa tres conceptos de reward para explicar el origen de la señal: el Process Reward Model (PRM, recompensa paso a paso), el Outcome Reward Model (ORM, que solo observa la respuesta final) y el RLHF / DPO tradicional (preferencia humana). El entrenamiento self-rewarding se acerca más al PRM, pues intenta estimar una puntuación después de cada paso de razonamiento del modelo, mientras que el Meta judge adopta la perspectiva del ORM para comparar el draft con la verified response. Weston recurre a ejemplos de benchmark como ``DeepSeek-R1, O1-R1'' para mostrar que, sin evaluación humana adicional, la autorrecompensa aún puede elevar el reasoning win rate (sobre todo en Alpaca Eval 2).

\begin{importantbox}{Comparación de modelos de recompensa}
\begin{itemize}
    \item PRM (recompensa de proceso): cada paso de la chain tiene señal, adecuado para ajustar el reasoning trace.
    \item ORM (recompensa de resultado): solo observa el juicio final, adecuado para casos como el meta judge, que solo necesita determinar cuál respuesta es mejor.
    \item RLHF / DPO: requiere preferencia humana, el costo de entrenamiento es alto, y aun así se sigue usando como baseline para la comparación.
\end{itemize}
\end{importantbox}

\begin{knowledgebox}{El valor de la recompensa verificable}
Weston señala la idea de ``verifiable reward'': si el modelo puede tratar su propia verified response como un ground truth aproximado, entonces se pueden usar estos data para train el reward model y reforzar la self-evaluation, en lugar de depender por completo de etiquetas humanas de bueno o malo.
\end{knowledgebox}

\subsection{Resumen de la sección}
El núcleo de la automejora consiste en reescribir la reward signal: dejar que el propio modelo determine el valor de la cadena de razonamiento, y conservar durante el entrenamiento tanto la process view como la outcome view. Los capítulos siguientes concretarán este objetivo de entrenamiento en un pipeline reutilizable.

\section{Pipeline de entrenamiento iterativo}
\subsection{Desglose de la estructura de Self-Rewarding}
El circuito cerrado de autoevaluación de recompensa incluye tres roles: el Task Generator se encarga de derivar tareas a partir de las semillas de Self-Instruct, el Solution Generator da la respuesta inicial con Chain-of-Thought y el Reward Model ejecuta la verificación cruzada y calcula la recompensa. Weston considera que las tareas producidas con Self-Instruct se acercan más a las necesidades reales (código, matemáticas, evaluación), por lo que este paso no debe limitarse a repetir el conjunto de datos de instrucciones ya existente.

\begin{figure}[H]
\centering
\includegraphics[width=0.8\textwidth]{frame-01-selfreward.jpg}
\caption{El presentador muestra en un capítulo anterior el diagrama de conceptos de self-improving, que ilustra el pipeline de tres roles.\protect\footnotemark}
\end{figure}
\footnotetext{Intervalo de tiempo en el video: 00:02:35--00:02:37.}

\subsection{Self-Instruct y validación cruzada}
Para mantener la amplitud de las tareas, el pipeline usa Self-Instruct para generar por lotes nuevos prompts. Después de cada ronda de respuestas, el modelo genera una verified response que se empareja con el borrador inicial. Si el meta judge determina que ambas difieren, se usa la verified response para actualizar el modelo de recompensa. Este proceso se asemeja a la validación conversacional del ``Self-Feeding Chatbot'', pero se enfoca más en la cadena de razonamiento que en la simple respuesta conversacional.

\begin{importantbox}{La sinergia entre Self-Instruct y la verified response}
\begin{itemize}
    \item Self-Instruct: a partir de un número reducido de seed prompts, duplica automáticamente tareas de matemáticas, código y evaluación.
    \item Verified response: el modelo determina si se equivocó; aunque el borrador original ya contenga la respuesta escrita, de todos modos debe generar una versión nueva y verificada.
\end{itemize}
\end{importantbox}

\begin{figure}[H]
\centering
\includegraphics[width=0.78\textwidth]{frame-02-iterative.jpg}
\caption{Ilustración del pipeline de Self-Rewarding, con énfasis en la actualización iterativa de la verified response.\protect\footnotemark}
\end{figure}
\footnotetext{Intervalo de tiempo en el video: 00:36:00--00:36:02.}

\begin{figure}[H]
\centering
\includegraphics[width=0.75\textwidth]{frame-03-selfinstr.jpg}
\caption{Flujo de generación de tareas y validación de respuestas dentro de Self-Instruct.\protect\footnotemark}
\end{figure}
\footnotetext{Intervalo de tiempo en el video: 00:44:40--00:44:42.}

\subsection{Datos de entrenamiento y estrategia de fine-tuning}
Weston menciona repetidamente en la clase el ``Self-Feeding Chatbot'', los ``SelfT evaluators'' y las ``verified response iterations'', componentes que forman el eje principal de los datos de entrenamiento. En cada ronda de self-rewarding se agrega una generación alr (instruction, question, answer, reasoning) y, después, bajo la indicación del meta judge, se filtran las respuestas de mayor calidad. Vale la pena notar que, en los experimentos reales, la coherencia entre el modelo entrenado con tres rondas de self-rewarding y la verified response mejoró considerablemente, tal como dijo Weston: ``\textit{I presently love the simplicity of it that we didn't have to do this large amount of Chain of Thought where it sort of overthinks itself}''.

\begin{knowledgebox}{Estrategia de muestreo de Self-Instruct}
Self-Instruct parte de un número reducido de seeds hechas a mano y genera, de forma cuasi supervisada, tripletas de instruction → input → output. Weston recomienda cubrir en cada ronda cuatro tipos de tareas: lenguaje, matemáticas, programación y evaluación lógica, para evitar que el conjunto de datos de entrenamiento se concentre en un solo tipo de razonamiento.
\end{knowledgebox}

\subsection{Resumen de la sección}
El pipeline de entrenamiento de Self-Rewarding conecta al Task Generator, al Solution Generator y al Reward Model en un solo circuito cerrado. Self-Instruct garantiza la diversidad de las tareas, la Verified Response hace que el modelo de recompensa se concentre en los patrones de error, y Self-Feeding/el meta judge verifica además si el resultado vale la pena adoptarse.
\section{Meta evaluación y generalización a largo plazo}
\subsection{Detalles del entrenamiento del Meta-Judge}
El meta judge surge porque los modelos self-rewarding llegan a un plateau después de few iterations: Weston lo describe de manera gráfica: ``self-rewarding models that I showed earlier only had the top part of the pipeline; introducing meta judge builds the bottom part, lifting the plateau'' (SRT 2756-2789). La entrada del meta judge es el par draft y verified response, y la salida indica cuál de los dos vale la pena conservar; el objetivo de entrenamiento ya no depende de preferencias humanas, sino que se refuerza según el win rate en benchmarks como Alpaca Eval 2.

\begin{figure}[H]
\centering
\includegraphics[width=0.77\textwidth]{frame-04-meta.jpg}
\caption{El capítulo del Meta Judge muestra la división de trabajo en tres pasos: draft, judge y verified.\protect\footnotemark}
\end{figure}
\footnotetext{Intervalo de tiempo en el video: 01:01:30--01:01:32.}

\subsection{Datos del Meta-Judge y SelfT Evaluators}
Weston menciona que ``SelfT evaluators son los que hicimos el año pasado en el paper de SelfT''\footnote{Traducción del hablante; en el original en inglés: ``SelfT evaluators''.}, un conjunto de datos que estructura el proceso de decisión sobre la verified response: incluye tres etapas, draft, verified y judge, y la etiqueta del judge proviene de tareas de razonamiento de escritorio y no de preferencias puramente humanas (SRT 2902-2904). El judge entrenado con los SelfT evaluators logra distinguir con más precisión las respuestas que ``realmente mejoran la calidad del reasoning'', con lo cual se eleva el win rate.

\begin{knowledgebox}{Lógica de diseño de los SelfT Evaluators}
Los SelfT Evaluators se entrenan con borradores que tienen chain-of-thought y respuestas ya verificadas; el judge solo tiene que elegir la más confiable entre ambas. Un evaluator entrenado así logra, con poca supervisión humana o con datos autosupervisados, una capacidad de discriminación mayor que la de los modelos de reward tradicionales.
\end{knowledgebox}

\subsection{Interpretabilidad del razonamiento y el reto computacional}
Weston sostiene que ``transformer itself does reasoning by all these steps'', pero al mismo tiempo reconoce que la red en sí no es transparente; por eso es necesario recurrir al meta judge o a un reward verificable para rastrear el reasoning trace (SRT 3152-3204). También insiste en que el ``inference time compute for evaluation'' es el cuello de botella del self-improvement, razón por la cual el número de verified response no puede aumentarse a la ligera, pues de lo contrario el costo de inferencia se duplicaría (SRT 3236-3252).

\subsection{Resumen de la sección}
El meta judge y los SelfT evaluators dotan a los modelos self-rewarding de una capacidad de discriminación genuina: al comparar el draft con la verified response, se asegura que la señal de reward no se confunda con hallucination, y al mismo tiempo se nos recuerda que el cómputo de inferencia debe mantenerse bajo control.

\section{Métricas experimentales clave}
\subsection{Observaciones de Iteration y win rate}
Weston divide el entrenamiento de self-rewarding en varias iteraciones; cada ronda combina tareas de Self-Instruct, verified response y judge. Los experimentos muestran que la primera ronda de entrenamiento ya le permite al modelo obtener un win rate considerable en Alpaca Eval 2, y que al agregar el meta judge el win rate mejora aún más, evitando el plateau que originalmente aparecía en la tercera ronda (SRT 1908-2027, 2784-2790).

\begin{table}[H]
\centering
\begin{tabular}{p{4cm}p{4cm}p{7cm}}
\toprule
Etapa de iteración & Estimación del tamaño de muestra & Observación \\
\midrule
Seed Iteration & Pocos verified prompts & El modelo domina rápido el tipo de tarea, pero la calidad del reasoning aún depende del judge. \\
Self-Rewarding Iteration 1 & Self-Instruct + verified response & El win rate mejora inicialmente en Alpaca Eval 2; el modelo aprende a autoevaluarse. \\
Self-Rewarding Iteration 2 & Tareas nuevas + cross-check & Se verifican los errores de manera repetida y el reward se desplaza hacia la dirección correcta. \\
Meta-Judge Iteration & Se agrega el judge de SelfT evals & Se resuelve el plateau, la confiabilidad del verified response se vuelve cuantificable y el win rate sube de manera constante. \\
\bottomrule
\end{tabular}
\caption{Observaciones de iteration y win rate presentadas en la clase (qualitative).}
\end{table}

\subsection{SelfT evaluators y otras fuentes de datos}
Además de Alpaca Eval 2, Weston también menciona SelfT, Self-Feeding Chatbot y generic instruction data (matemáticas/código/juicio/conversación). Estas fuentes de datos se complementan entre sí: SelfT aporta judge supervision, Self-Feeding aporta conversational feedback y Self-Instruct aporta new tasks, y forman los cuatro puntos de apoyo necesarios para el «autoentrenamiento del modelo».

\begin{importantbox}{Los cuatro puntos de las fuentes de datos}
\begin{itemize}
    \item Self-Instruct: genera instruction → input → output.
    \item Self-Feeding Chatbot: genera verified response mediante la conversación.
    \item SelfT evaluators: aporta etiquetas simuladas de judge.
    \item Alpaca Eval 2: observación del win rate en el benchmark, que se usa para guiar el modelo de reward.
\end{itemize}
\end{importantbox}

\subsection{Resumen de la sección}
Las distintas iteraciones intensifican progresivamente tanto las fuentes de datos como el judge: Self-Instruct se encarga de las tareas, Self-Feeding aporta la verificación, SelfT entrena al judge y Alpaca Eval 2 aporta el win rate final; los cuatro elementos trabajan juntos para mejorar de manera constante el desempeño en reasoning.

\section{Casos prácticos y diagnóstico de problemas}
\subsection{Plantilla de Self-Feeding Chatbot}
Weston vuelve a revisar la práctica de Self-Feeding Chatbot: después de que el modelo dialoga con el usuario, convierte el draft inicial en una verified response, y luego incorpora las respuestas de alta calidad al conjunto de entrenamiento para apoyar el self-rewarding (SRT 774-884). Este paso enfatiza ``rewarding itself following instructions''; el modelo trata la verified response como otro tipo de label y la refuerza en future iterations.

\begin{table}[H]
\centering
\begin{tabular}{p{4cm}p{10cm}}
\toprule
Etapa & Descripción \\
\midrule
User Dialogue & El modelo interactúa con el usuario y registra la entrada, la salida y el contexto, lo que forma una señal de supervisión inmediata. \\
Model Draft & El modelo genera un draft y produce un Chain-of-Thought, que sirve como base del juicio. \\
Self-Verification & El modelo vuelve a revisar el draft, genera una verified response y enfatiza la consistencia lógica. \\
Verified Response & Esta respuesta se trata como una muestra de alta calidad para que el reward model aprenda de ella. \\
Training Update & La verified response se incorpora al conjunto de entrenamiento para seguir optimizando el reward / meta judge. \\
\bottomrule
\end{tabular}
\caption{Etapas clave y flujo de datos en la plantilla de Self-Feeding.}
\end{table}

\subsection{Patrones de error y alertas}
En las iteraciones de self-rewarding son comunes tres tipos de problemas: 1) overthinking del modelo, es decir, entre más piensa, más se aleja de la respuesta; 2) hallucination combinada con exceso de confianza; 3) meta judge drift, por ejemplo cuando un calibration drift interno del judge hace que incluso una verified response se juzgue erróneamente como mala. Weston también menciona ``we can use judge to cross check itself and see that it was wrong'' (SRT 1188-1190), lo que indica que necesitamos mantener un check-list confiable para detectar anomalías.

\begin{table}[H]
\centering
\begin{tabular}{p{4cm}p{5cm}p{5cm}}
\toprule
Patrón de error & Manifestación típica & Estrategia de respuesta \\
\midrule
Overthinking & El Chain-of-Thought se alarga gradualmente sin que la respuesta mejore & Limitar el reasoning depth e introducir un meta judge para corregir el rumbo \\
Hallucination + Confidence & El modelo se muestra seguro pero la respuesta se desvincula de los hechos & Usar múltiples evaluator y aumentar la autoverificación \\
Meta judge drift & Incluso una verified response se juzga como mala & Recalibrar periódicamente los datos del SelfT judge \\
\bottomrule
\end{tabular}
\caption{Alertas clave que deben monitorearse en la práctica de Self-Rewarding.}
\end{table}

\subsection{Resumen de la sección}
Self-Feeding Chatbot ofrece una plantilla de verificación en línea, y la tabla de patrones de error nos recuerda que es necesario recalibrar mediante el meta judge, limitar el reasoning depth y reforzar el monitoreo de hallucination para mantener el self-improvement en la dirección correcta.

\section{Lista de verificación del proceso y validación}
\subsection{Preparación previa al entrenamiento}
Antes de entrar al entrenamiento self-rewarding, las recomendaciones de Weston incluyen: elegir las seeds de Self-Instruct, preparar los evaluators de SelfT, configurar el inference compute budget y establecer la frecuencia de cross-check. Estos pasos permiten que cada iteración genere de manera controlada una verified response, en la que «las tareas producidas por Self-Instruct deben cubrir matemáticas, programación y juicio» (SRT 674-676).

\begin{itemize}
    \item Establecer la seed de Self-Instruct (Math / Code / Reasoning / Dialogue) y ampliarla hasta obtener diverse prompts.
    \item Preparar el pipeline de verified response y definir los criterios de etiquetado de cross-check.
    \item Preentrenar o hacer fine-tuning del SelfT evaluator, para garantizar que el judge pueda distinguir entre draft y verified.
    \item Planificar el inference-time compute, para evitar que la salida de cross-check se convierta en un compute bottleneck (SRT 3236-3252).
    \item Respaldar los datos de conversación de Self-Feeding, para usarlos en la supervisión posterior.
\end{itemize}

\subsection{Procesamiento posterior a la evaluación}
Al terminar cada ronda de entrenamiento, es necesario usar los results para calibrar el meta judge, calcular el win rate y registrar si la verified response es consistente entre iteration. Weston menciona en particular el proceso de ``kept verifying itself and iterating'' (SRT 1908-2027), y recomienda extraer la verified response para compararla con Alpaca Eval 2, y volver a puntuarla con el judge, para asegurar que se rompa el plateau.

\begin{itemize}
    \item Hacer que el meta judge vuelva a puntuar draft y verified response, para capturar el drift en tiempo real.
    \item Comparar la verified response más reciente con Alpaca Eval 2 / Alpaca Eval 2 baseline.
    \item Reunir los error logs, prestando especial atención a las muestras de hallucination y overthinking.
    \item Ajustar el peso del loss del reward model y la frecuencia de cross-check según la retroalimentación del judge.
    \item Marcar las verified response de alta calidad como future seed, para seguir ampliando el corpus de Self-Instruct.
\end{itemize}

\subsection{Resumen de la sección}
La preparación previa al entrenamiento se encarga de garantizar que las distintas tareas y el judge estén listos, mientras que el procesamiento posterior a la evaluación se encarga de la corrección y la síntesis: solo si la verified response se usa tanto para el entrenamiento como para la evaluación, el self-rewarding puede seguir impulsando la calidad del reasoning.

\section{Inspiración para la investigación y desafíos de despliegue}
\subsection{La fusión entre el autosupervisado y la validación en producción}
El ponente reiteró el flujo de trabajo de «Self-Feeding Chatbot»: el modelo se etiqueta a sí mismo después de la conversación, genera una verified response y luego incorpora las respuestas de alta calidad al entrenamiento, con lo que logra online self-improvement (SRT 774-884). La diferencia entre este mecanismo y el entrenamiento self-rewarding es que la señal de supervisión proviene de interacciones reales con usuarios, y no de tasks generadas de manera pura dentro del conjunto de entrenamiento.

\begin{warningbox}{Riesgo de robustez en el despliegue}
La hallucination sigue siendo una threat: si el judge considera razonable una respuesta y la hallucination resulta ser autoconsistente, todo el proceso de self-improvement puede entrar en un ciclo de amplificación de una «falacia lógica reforzada». Es necesario introducir un evaluator externo o evidencia multimodal para acotar este tipo de errores.
\end{warningbox}

\subsection{Problemas de frontera y direcciones futuras}
Weston mencionó varias direcciones futuras, entre ellas inference time compute for evaluation, interpretability y sistemas «self-aware». En concreto, dijo: «Reason by actually interacting with people in the world, internet or itself» (SRT 3282-3284), y señaló la relación entre el presupuesto de compute de la self evaluation y el reasoning understand (SRT 3236-3252). El autoconocimiento (self-aware), aunque sigue siendo algo impreciso, implica al menos que el modelo pueda buscar ayuda externa de manera activa cuando carece de un judge.

\begin{importantbox}{Tres prioridades del trabajo futuro}
\begin{itemize}
    \item El inference-time compute budget del razonamiento y la evaluación.
    \item Un pipeline de interpretation transparente, que permita a las personas hacer trace del razonamiento.
    \item Sistemas self-aware: el modelo sabe buscar un external agent cuando la uncertainty es alta.
\end{itemize}
\end{importantbox}

\subsection{Resumen de la sección}
La investigación debe equilibrar el self-training offline con el verified feedback online, y buscar un balance entre hallucination, compute e interpretability. El objetivo final es construir un modelo que razone y a la vez sea capaz de autoevaluarse.

\section{Diseño y ejemplos de tareas de razonamiento}
\subsection{Ejemplo de razonamiento en varios pasos}
Para que el pipeline de self-rewarding aprenda razonamiento en varios pasos, Weston propone dividir el Chain-of-Thought en tres pasos: «decompose → verify → finalize». Por ejemplo:
\begin{itemize}
    \item Prompt: «Analiza un problema algebraico de 4 pasos, mostrando cada paso con variables intermedias».
    \item Draft: enumera la cadena de razonamiento paso a paso y da una respuesta inicial.
    \item Verification: el modelo vuelve a revisar el draft y, cuando no puede verificarlo, genera una idea alternativa.
\end{itemize}
\textit{``You can make it kind of cross check itself and see that it was wrong and write a final verified response.''}(SRT 1188-1190)Este proceso es la aplicación típica de multi-step reasoning + meta judge.

\begin{table}[H]
\centering
\begin{tabular}{p{4cm}p{4cm}p{6cm}}
\toprule
Tipo de tarea & Prompt de ejemplo & Métrica de verificación \\
\midrule
Razonamiento algebraico & «Descompón x+2y=7, x-y=1» & Integridad del Chain-of-Thought, consistencia de la respuesta final \\
Análisis causal & «Describe cómo un agente de RL reflexiona sobre el reward hacking» & Si señala el punto de desviación en los pasos de la operación \\
Juicio estadístico & «Compara la precisión y la confianza de dos modelos» & Consistencia del meta judge al determinar el winner \\
\bottomrule
\end{tabular}
\caption{Tareas típicas de multi-step reasoning en Self-Instruct.}
\end{table}

\begin{knowledgebox}{Control de calidad del razonamiento en varios pasos}
Cada cadena de razonamiento debe llevar consigo una verified response, de modo que el modelo de reward pueda tratar los errores como señal de supervisión, en lugar de premiar solo la extensión o la fluidez. Tanto una reasoning chain demasiado corta como una chain con overthinking deben perder puntos en el cross-check del judge.
\end{knowledgebox}

\subsection{Ejemplo de programación y verificación}
En las tareas de programación, el scene de Self-Instruct suele exigir que el modelo entregue código completo, explicación y pruebas. El pipeline de self-rewarding hace que el modelo, después de generar el código, vuelva a «ejecutar» su propio razonamiento (por ejemplo, simulando casos de prueba) y luego dé una verified response que explique por qué ese código cumple los requisitos. Esa verified response de paso le enseña al modelo a hacer self-evaluate de su output.

\begin{itemize}
    \item Prompt: «Escribe una función de Python que reciba una lista y devuelva los elementos únicos».
    \item Draft: devuelve la función más una breve explicación.
    \item Meta judge: usa el SelfT evaluator para comparar el draft con el verified, y determina cuál explicación es más clara y cuál código es más seguro.
\end{itemize}

\subsection{Escenario simulado del meta judge}
En el deployment real, el meta judge desempeña el papel del «último guardrail». Por ejemplo, una vez que el draft y el verified están terminados, el judge calcula la diferencia de alignment entre ambos y da un tick con alta confianza. Este check a veces requiere introducir un external agent (como un evaluador humano o de terceros), sobre todo en los casos en que el judge tiene confianza pero hallucinates (SRT 2500-2509).

\subsection{Resumen de la sección}
Mediante multi-step reasoning, tareas de programación y escenarios simulados de meta judge, se puede tratar el entrenamiento de self-rewarding como un diseño de alto nivel, que va encadenando prompt → draft → verified → judge hasta formar una plantilla de entrenamiento repetible.

\section{Elementos visuales y gestión de recursos}
\subsection{Fotogramas clave y marcas de tiempo}
Esta clase se apoya en la imagen de portada y en 4 fotogramas clave para reforzar la evidencia visual. Al no haber slides, estos fotogramas muestran conceptos centrales como self-rewarding, self-instruct y meta-judge. La tabla adjunta enumera el nombre del fotograma, el intervalo de tiempo y la descripción textual, para facilitar su consulta posterior.

\begin{table}[H]
\centering
\begin{tabular}{p{3cm}p{3cm}p{7cm}}
\toprule
Archivo del fotograma & Tiempo & Descripción \\
\midrule
frame-01-selfreward.jpg & 00:02:35--00:02:37 & Diagrama conceptual del self-improvement pipeline \\
frame-02-iterative.jpg & 00:36:00--00:36:02 & Bucle iterativo de Self-Rewarding \\
frame-03-selfinstr.jpg & 00:44:40--00:44:42 & Tareas de Self-Instruct y verified response \\
frame-04-meta.jpg & 01:01:30--01:01:32 & Meta Judge compara draft y verified response \\
\bottomrule
\end{tabular}
\caption{Fotogramas clave y marcas de tiempo usados en esta clase.}
\end{table}

\subsection{Subtítulos y apoyo para la redacción}
A partir de \texttt{lecture11.en.srt} depuramos \texttt{subs\_clean.txt}, extrayendo oraciones temáticas en bloques de 2500 líneas y dividiendo después las secciones según la lógica didáctica. La versión clean del texto reduce las repeticiones (como «days I'm using AI...»), lo que aporta material estable para ampliar el contenido de cada sección.

\subsection{Resumen de la sección}
Los recursos visuales y la depuración de subtítulos constituyen el doble sustento del material complementario: los primeros satisfacen el requisito del usuario de frame + tiempo; los segundos dan un fundamento a la redacción de las secciones y mantienen la consistencia de los detalles técnicos y del orden.

\section{Temas adicionales: guía y calibración}
\subsection{Ingeniería de prompts de Self-Instruct}
Para que los prompts generados por Self-Instruct cubran a la vez matemáticas, código, razonamiento lógico y conversación, Weston recomienda escribir explícitamente en el prompt ``show your chain of thought'', ``highlight where you doubt'' y ``produce a verification statement''. Entre las prácticas concretas están:
\begin{itemize}
    \item Usar la plantilla: ``Step-by-step, explain ..., then verify with a judgement block.''.
    \item Agregar al final del prompt ``If you are unsure, state your uncertainty and ask for a cross-check''.
    \item Fijar de antemano el formato de la verified response (reasoning + final verdict), para que el entrenamiento del reward model pueda emparejarla con más facilidad.
\end{itemize}
\mbox{}

\subsection{Calibration y normalización de la recompensa}
Entrenar el meta judge requiere calibration, para poder juzgar con justicia entre la verified response y el draft. Weston nos recuerda que ``self evaluation is only going to help self-improve the model if inference compute for evaluation is controlled'' (SRT 3236-3252). Entre las estrategias habituales están reward clipping, temperature scaling, cross-check rate scheduling y el reentrenamiento del judge.

\begin{knowledgebox}{Calibration best practices}
Aplicar temperature scaling a los logits del judge, o elevar la frecuencia de cross-check cuando la verified response presenta baja confidence; aplicar clipping y normalization a los valores de reward, para evitar que un solo outlier haga que el self-rewarding se descontrole por completo.
\end{knowledgebox}

\subsection{Resumen de la sección}
El diseño de prompts y la calibration son condiciones necesarias para el éxito del self-rewarding: el primero garantiza la cobertura de las tareas sin perder el rumbo, y la segunda hace que el judge sea más estable durante el cross-check y no se desvíe por exceso de confianza ante una hallucination.
\section{Open Questions y direcciones a largo plazo}
\subsection{Interpretabilidad y regulación}
Weston considera que ``the transformer itself does reasoning'', pero también nos advierte que ``interpreting the neural network isn't easy'' (SRT 3204). Por ello, una de las direcciones a largo plazo es construir un pipeline de razonamiento trazable, de manera que el verified response, además del final answer, también pueda producir un trace para uso regulatorio.

\begin{table}[H]
\centering
\begin{tabular}{p{5cm}p{10cm}}
\toprule
Problema & Descripción \\
\midrule
Brecha de interpretabilidad & ¿Cómo extraer un reasoning trace del verified response para que un humano lo revise? \\
Atribución de responsabilidad & Cuando el meta judge se equivoca, ¿cómo rastrear si la falla proviene del reward model, del judge o de una operación incorrecta del prompt? \\
Señales multiagente & En un escenario con múltiples agentes (bot ↔ judge ↔ user), ¿cómo garantizar que la reward signal no sea objeto de adversarial exploitation? \\
\bottomrule
\end{tabular}
\caption{Open Questions para uso posterior en alignment y auditoría.}
\end{table}

\subsection{Interacción con sistemas externos}
Los modelos que se automejoran en el futuro necesitan ``Reason by actually interacting with people in the world, internet or itself'' (SRT 3282-3284). Esto incluye:
\begin{itemize}
    \item Coordinarse con otros agentes (share verified responses, cross judge);
    \item Recabar additional context mediante API / internet, para después autoverificarse dentro de Self-Instruct;
    \item Buscar automáticamente la intervención de un agente externo al detectar una hallucination (self-aware behaviour).
\end{itemize}

\subsection{Resumen de la sección}
Open Questions ofrece una dirección para el alignment y la regulación; la interacción con el mundo exterior es el siguiente paso para construir un sistema self-improving más autónomo, más autoconsciente y también más seguro.

\section{Apéndice: recopilación de datos y procesamiento de imágenes}
\subsection{Descarga de video y extracción de fotogramas}
Para obtener una portada nítida y los key frames, descargamos el video original con precisión de 720p mediante \texttt{yt-dlp -f 398+251 -o lecture11.webm} (enlace del video \texttt{https://www.youtube.com/watch?v=\_MNlLhU33H0}); luego, en los cuatro instantes 00:02:35, 00:36:00, 00:44:40 y 01:01:30, invocamos \texttt{ffmpeg -ss ... -frames:v 1 -update 1} para generar sucesivamente frame-01–frame-04. La tabla siguiente enumera los comandos clave y su propósito.

\begin{table}[H]
\centering
\begin{tabular}{p{6cm}p{9cm}}
\toprule
Comando & Descripción \\
\midrule
\texttt{yt-dlp -f 398+251 -o lecture11.webm https://www.youtube.com/watch?v=\_MNlLhU33H0} & Descarga el video en 720p+audio, y proporciona el material para las capturas y la portada posteriores. \\
\texttt{ffmpeg -ss 00:02:35 -i lecture11.webm -frames:v 1 -update 1 frame-01-selfreward.jpg} & Captura el key frame de la sección de self-rewarding. \\
\texttt{ffmpeg -ss 00:36:00 -i lecture11.webm -frames:v 1 -update 1 frame-02-iterative.jpg} & Captura la ilustración del iterative pipeline. \\
\texttt{ffmpeg -ss 00:44:40 -i lecture11.webm -frames:v 1 -update 1 frame-03-selfinstr.jpg} & Fotograma de ejemplo de Self-Instruct. \\
\texttt{ffmpeg -ss 01:01:30 -i lecture11.webm -frames:v 1 -update 1 frame-04-meta.jpg} & Imagen de Meta judge. \\
\bottomrule
\end{tabular}
\caption{Comandos clave de recopilación e imágenes.}
\end{table}

\subsection{Depuración de subtítulos y esquema}
Depuramos el SRT mediante Python: eliminamos los párrafos duplicados, conservamos las marcas de tiempo y extraemos la oración temática cada 2,500 líneas por bloque para apoyar la escritura. El núcleo del script divide los bloques con \texttt{re.split(r'\\n\\n+')}, extrae la marca de tiempo de la primera línea y usa \texttt{prev} para eliminar duplicados y conservar lo esencial. El archivo \texttt{subs\_clean.txt} resultante sirve, en la redacción de las notas, como fuente de la chronology; contiene palabras clave de párrafo como \texttt{Self-Rewarding} y \texttt{Meta Judge}, lo que facilita clasificarlo según la lógica didáctica.

\subsection{Resumen de la sección}
El proceso de descarga, captura de pantalla y procesamiento de subtítulos que documenta el apéndice constituye el respaldo de la escritura posterior, y garantiza que cada imagen y cada fragmento de texto puedan rastrearse hasta un instante y un script específicos.
\section{Apéndice: terminología y abreviaturas}
\subsection{Explicación de los términos principales}
\begin{table}[H]
\centering
\begin{tabular}{p{4cm}p{11cm}}
\toprule
Término & Explicación \\
\midrule
Chain-of-Thought & Hace que el modelo enumere los pasos de razonamiento intermedios, lo cual proporciona la base para el self-check. \\
Self-Rewarding & El modelo genera tareas por sí mismo y utiliza el reward para mejorar el reasoning, sin depender de labels manuales. \\
Self-Instruct & A partir de unos cuantos seeds, genera automáticamente instruction→input→output, con lo que amplía el corpus de entrenamiento. \\
Verified Response & La respuesta final que el modelo genera después de autoverificarse, la cual sirve como señal confiable para el reward. \\
Meta Judge & Un evaluator que compara el draft con el verified y decide cuál es más confiable. \\
SelfT Evaluators & Un conjunto de datos de judge entrenado usando el verified response como etiqueta. \\
\bottomrule
\end{tabular}
\caption{Varios términos mencionados repetidamente en las notas de esta clase.}
\end{table}

\subsection{Abreviaturas relacionadas}
\begin{itemize}
    \item PRM (Process Reward Model): recompensa paso a paso, que se usa para incentivar cada paso del reasoning.
    \item ORM (Outcome Reward Model): solo considera la respuesta final; es el enfoque típico del Meta judge.
    \item RLHF (Reinforcement Learning from Human Feedback): entrenamiento tradicional de reward que incluye retroalimentación humana.
    \item DPO (Direct Preference Optimization): una estrategia de actualización de parámetros que no usa RL y también aplica preference learning.
\end{itemize}

\subsection{Resumen de la sección}
El apéndice de terminología y abreviaturas ayuda a repasar las palabras clave de la clase, y facilita ubicar rápidamente la relación entre Self-Rewarding, Meta Judge y Self-Instruct.

\section{Guía de monitoreo y evaluación}
\subsection{Métricas clave}
El entrenamiento de automejora debe monitorear continuamente las siguientes métricas, para poder intervenir con rapidez cuando el win rate llegue a una meseta, la hallucination aumente o el inference compute se dispare.

\begin{table}[H]
\centering
\begin{tabular}{p{5cm}p{10cm}}
\toprule
Métrica & Descripción \\
\midrule
Win rate delta & Con Alpaca Eval 2 / DeepSeek benchmarks como referencia, observar el cambio en el win rate entre el verified response y el draft. \\
Reward variance & Si la desviación estándar del reward score en el self-rewarding aumenta de forma abrupta, indica que el modelo podría estar en overthinking o hallucinating. \\
Hallucination frequency & Usar el evaluador SelfT para etiquetar la tasa de muestras con hallucination verificada, y reducir la cross-check depth cuando se supere el umbral. \\
Inference compute & Evaluar el compute adicional que genera el verified response (número de tokens × número de cross-checks), para evitar que el costo de la evaluation exceda el presupuesto. \\
\bottomrule
\end{tabular}
\caption{Métricas centrales para monitorear el entrenamiento self-rewarding.}
\end{table}

\begin{importantbox}{Recomendación sobre la interacción entre métricas}
Cuando el win rate se estanque, revisar primero el reward variance y la hallucination frequency: si la proporción de hallucination aumenta, retroceder a las Self-Instruct seeds; si el reward variance sube, depender de la meta judge calibration para converger a tiempo.
\end{importantbox}

\subsection{Alertas automatizadas y retroalimentación}
Una vez establecidos los umbrales para las métricas correspondientes, se pueden usar automated scripts para generar alert: por ejemplo, pausar la iteración cuando el reward variance supere 0.3, o activar una revisión human-in-the-loop cuando la hallucination frequency sea > 5\% durante tres veces consecutivas. Esto puede combinarse con el flujo de Self-Feeding, enviando el alert al pipeline dashboard, de modo que el equipo de operación pueda cerrar el ciclo con rapidez dentro del \texttt{self-improvement}.

\begin{knowledgebox}{Plantilla de estrategia de alertas}
1. Definir la frecuencia: calcular las métricas inmediatamente al terminar cada iteration.\newline
2. Respuesta escalonada: el warning activa un retry prompt; el alert activa un verified response adicional.\newline
3. Archivar los registros: las muestras de hallucination junto con la trayectoria del reward se escriben en \texttt{audit.csv} para el meta judge retraining futuro.
\end{knowledgebox}

\subsection{Resumen de la sección}
Con este conjunto de métricas y alertas, el equipo puede dar seguimiento de manera transparente al estado del self-rewarding durante el entrenamiento real, y en cuanto aparezca un drift, puede recalibrar con el meta judge o repetir el self-instruct prompt.
\section{Estrategia de despliegue y gobernanza}
\subsection{Estrategia de salida del modelo}
Weston enfatiza que la salida de self-rewarding no debe exponer directamente la raw verified response: es necesario pasar primero por una capa de alignment (meta judge + human checker) que la confirme antes de producir la salida. La estrategia de despliegue suele incluir:
\begin{itemize}
    \item Deployment filters: solo puede ingresar a la downstream API la verified response que pase la verificación de consistencia con el judge.
    \item Conservative fallback: cuando el judge confidence esté por debajo del umbral, se degrada automáticamente al modelo baseline RLHF.
    \item Update schedule: se genera una versión self-rewarding a diario o semanalmente, y se prueba retrospectivamente la hallucination con el offline dataset.
\end{itemize}

\begin{warningbox}{Advertencia de riesgo de despliegue}
Si el modelo self-rewarding presenta hallucination de manera persistente en producción, aunque la verified response haya sido aceptada por el judge, de todos modos puede producir misleading explanations en la conversación. Se necesita un multi-agent cross-check o una nightly calibration para evitar el reward model drift.
\end{warningbox}

\subsection{Gobernanza y auditoría}
A nivel de gobernanza se debe registrar: 1) el judge log de cada verified response; 2) el reward score y el compute cost; 3) la versión del prompt template y de la Self-Instruct seed. Con ayuda de estos log, se puede aportar transparencia para future audits y sustentar la pregunta `Why did the model take this reasoning path?`

\begin{table}[H]
\centering
\begin{tabular}{p{5cm}p{10cm}}
\toprule
Dimensión de gobernanza & Práctica \\
\midrule
Audit trail & Guardar la verified response, los resultados del judge y el reward score para reproducirlos posteriormente. \\
Prompt provenance & Marcar cada Self-Instruct prompt con un seed ID y un generation timestamp. \\
Compute budget governance & Establecer un máximo de tokens para el cross-check; si se excede, se dispara una alert. \\
\bottomrule
\end{tabular}
\caption{La capa de gobernanza hace que el self-improvement sea auditable.}
\end{table}

\subsection{Resumen de la sección}
Las políticas de filtrado previas al despliegue y los registros de gobernanza aseguran que la verified response no genere confusión directamente, y permiten que el equipo pueda responder en cualquier momento ``How did this reasoning chain earn its reward''.

\section{Síntesis y ampliación}
\begin{table}[H]
\centering
\begin{tabular}{p{4.5cm}p{11cm}}
\toprule
Tema & Puntos clave de la clase \\
\midrule
Arquitectura de autorrecompensa & El Task Generator genera tareas diversas, el Solution Generator produce Chain-of-Thought, el Reward Model se autoevalúa mediante verified response. \\
Self-Instruct + Verified Response & Self-Instruct asegura la amplitud de las tareas, verified response aporta muestras de cross-check y evita el encadenamiento excesivo. \\
Meta-Judge y evaluación & Meta-Judge usa SelfT evaluators para simular al juez, logra distinguir entre draft y verified response y supera de forma efectiva el plateau. \\
Retos de despliegue & el inference-time compute debe ser controlable, la alucinación requiere multi-evaluator y un sistema self-aware puede solicitar la intervención humana de forma más razonable. \\
\bottomrule
\end{tabular}
\caption{Hilo conductor central del self-improving pipeline visto en esta clase.}
\end{table}

\begin{table}[H]
\centering
\begin{tabular}{p{5cm}p{10cm}}
\toprule
Dirección de investigación & Descripción \\
\midrule
Recompensa verificable & Construir verified response para que el modelo de recompensa mantenga su confiabilidad sin necesidad de anotación humana. \\
SelfT + Meta Judge & Usar SelfT evaluators para entrenar al meta judge, recuperando de forma automática un juicio de alta calidad a partir del draft y el verified response. \\
Interacción en línea & Mediante el mecanismo Self-Feeding se enlazan la interacción real, el verified response y los datos de instrucción, formando un ciclo de aprendizaje continuo. \\
\bottomrule
\end{tabular}
\caption{Hoja de ruta de investigación a la que conviene dar seguimiento continuo.}
\end{table}

\subsection{Lecturas adicionales}
\begin{itemize}
    \item El artículo de OpenAI «Self-Feeding Chatbot», que muestra la iteración de verified response en el diálogo.
    \item La investigación de SelfT Evaluators / Self-Rewarding de Weston y colegas, que explica en detalle el entrenamiento del juez.
    \item El benchmark Alpaca Eval 2, que ofrece una referencia de win rate junto con la verificación mediante meta judge.
    \item Artículos comparativos entre Chain-of-Thought y RLHF/DPO, útiles para decidir cuándo usar self-rewarding.
\end{itemize}

\subsection{Resumen de la sección}
Esta clase de Jason Weston nos muestra que la capacidad de razonamiento es una cadena que necesita tres pasos: plantear la pregunta, verificar la respuesta y evaluar el resultado. Solo cuando el modelo aprende por sí mismo a partir del verified response y vuelve a confirmarlo bajo la guía del meta judge, el entrenamiento de reasoning puede avanzar al siguiente nivel.

\end{document}
